\documentclass[11pt]{article}
% =========================================================
%  學生講義 共用前言（以 XeLaTeX 編譯）
%  需要套件：xeCJK, tcolorbox（其餘多在 BasicTeX 內）
% =========================================================
\usepackage[a4paper,margin=2.3cm]{geometry}
\usepackage{amsmath,amssymb}
\usepackage{xcolor}
\usepackage{graphicx}

% ---- 中文字型（macOS 系統字型）----
\usepackage{xeCJK}
\setCJKmainfont{Songti TC}      % 內文：宋體
\setCJKsansfont{Heiti TC}       % 標題/強調：黑體
\xeCJKsetup{AutoFakeBold=2.2, AutoFakeSlant=0.2}

% ---- 版面 ----
\setlength{\parindent}{0pt}
\setlength{\parskip}{0.55em}
\linespread{1.15}
\renewcommand{\baselinestretch}{1.15}

% ---- 顏色 ----
\definecolor{cBlue}{HTML}{1f6feb}
\definecolor{cGray}{HTML}{6b7280}
\definecolor{cGrayBg}{HTML}{f3f4f6}
\definecolor{cPurp}{HTML}{7a3ff2}
\definecolor{cPurpBg}{HTML}{f5f1ff}

% ---- 標註框 ----
\usepackage[most]{tcolorbox}

% 就地補充新概念（灰底）：\begin{補充框}{標題}
\newtcolorbox{補充框}[1]{
  enhanced, breakable, arc=2pt, boxrule=0.6pt,
  colback=cGrayBg, colframe=cGray,
  left=9pt,right=9pt,top=7pt,bottom=7pt,
  before upper={\textbf{\sffamily #1}\par\vspace{3pt}},
}

% 延伸 / 開放問題（紫色左邊條）：\begin{延伸框}{標題}
\newtcolorbox{延伸框}[1]{
  enhanced, breakable, arc=2pt, boxrule=0pt,
  colback=cPurpBg, colframe=cPurpBg,
  borderline west={3pt}{0pt}{cPurp},
  left=10pt,right=9pt,top=7pt,bottom=7pt,
  before upper={\textbf{\sffamily 延伸閱讀｜#1}\par\vspace{3pt}},
}

% ---- 圖：置中、底下小字說明 ----
% \fig{檔名}{寬度}{說明}
\newcommand{\fig}[3]{%
  \begin{center}
  \includegraphics[width=#2\linewidth]{#1}\\[2pt]
  {\small\color{cGray} #3}
  \end{center}}

% ---- 章節標題用黑體 ----
\newcommand{\h}[1]{\sffamily #1}


% 讓羅馬數字 Ⅰ（U+2160）由中文字型（Noto CJK TC，含此字符）排印，避免 Latin sans 缺字
\xeCJKDeclareCharClass{CJK}{"2160}

\begin{document}

\begin{center}
{\sffamily\Huge\bfseries 選物I-3 平面運動}\\[3pt]
{\sffamily\large 普通高中 選修物理Ⅰ（力學一）・學生講義}
\end{center}
\vspace{0.6em}

直線運動只沿一條直線進行，位置用一個數字 $x$ 就能描述。但多數實際運動發生在平面上：被打出的棒球、轉彎的車、飛過操場的紙飛機，都同時具有「往前」與「往下（或往上）」的運動。要描述平面上的運動，一個數字不夠，須用兩個分量。

本章的核心觀念只有一個：\textbf{運動可以拆成水平、鉛直兩個獨立的方向}，各自用一維公式處理，最後再用向量合起來。把這個觀念用在「拋出去的物體」上，就得到本章的主角——\textbf{拋體運動（projectile motion）}，其軌跡為一條拋物線。學會「分方向處理」之後，下一章把力也分成兩個分量解題，是同一方法的延伸。

\medskip
本章主線：

\begin{center}
向量分量與運動獨立性 $\rightarrow$ 平拋運動 $\rightarrow$ 斜向拋射 $\rightarrow$ 綜合分析 $\rightarrow$ 平面相對運動
\end{center}

\section*{\sffamily 3-1\quad 平面運動與向量分量}

\subsection*{\sffamily A.\ 從一條線到一個平面}

要描述平面上的運動，最方便的做法是建立一組 $x$–$y$ 直角座標，把每個運動學量都寫成\textbf{向量}，再拆成沿 $x$、沿 $y$ 的兩個\textbf{分量（component）}。位移、速度、加速度都有兩個分量：
$$\vec r=(x,\,y),\qquad \vec v=(v_x,\,v_y),\qquad \vec a=(a_x,\,a_y).$$
速度大小（速率）為
$$|\vec v|=\sqrt{v_x^2+v_y^2},\qquad \text{方向 } \tan\theta=\frac{v_y}{v_x}.$$
向量的合成與分解在數學中已學過，這裡是把它用在運動上。

\subsection*{\sffamily B.\ 本章的核心：運動的獨立性}

本章最重要的性質是：\textbf{水平方向的運動和鉛直方向的運動，彼此互不影響}。$x$ 方向如何運動，不會影響 $y$ 方向如何運動，反之亦然。因此可以把一個平面運動拆成兩個各自獨立的一維運動，分別用等速度、等加速度公式計算，最後再把結果合起來。這個性質稱為\textbf{運動的獨立性（independence of motions）}。

\begin{補充框}{為什麼水平和鉛直可以「分開算、互不影響」？}
這並非理所當然，而是一個須說明的物理事實，其根源在於\textbf{重力只有鉛直分量，在水平方向上不施力}。

\textbf{物理上：}拋體在空中（忽略空氣阻力）只受重力，而重力是純鉛直的，所以水平方向沒有任何力。沒有力，水平速度就不會改變，為等速度運動；而「往下掉」這件事，重力給每個物體同樣的加速度 $g$，與它水平方向的快慢無關。

\textbf{數學上：}加速度是向量，兩個分量各自獨立相加。拋體只受純鉛直的重力，故
$$a_x=0,\qquad a_y=-g.$$
於是水平方向 $a_x=0\Rightarrow v_x$ 不變（等速度）；鉛直方向 $a_y=-g\Rightarrow$ 等加速度。兩條方程式中，$x$ 的式子不出現 $y$，$y$ 的式子不出現 $x$，這就是「獨立」的數學意思。（可類比為兩台互不相通的碼錶：一台只記水平位置、一台只記鉛直高度，重力只去動鉛直那台。）

\textbf{注意：}「球往前飛得越快，就越能撐住、比較不容易掉下來」並不正確。水平的快慢與鉛直的下墜是兩回事。此性質的前提是\textbf{忽略空氣阻力}；空氣阻力的方向沿速度的反方向，會同時牽涉水平與鉛直分量，使兩方向不再獨立（見 3-4 的延伸閱讀）。
\end{補充框}

\fig{si3-independence}{0.62}{圖 3-1：把一顆球每隔相同時間拍一次的頻閃示意。水平方向每格間距相同（等速度）；鉛直方向每格落差越來越大（等加速度，自由落體）。兩者合起來即為一條拋物線。}

\section*{\sffamily 3-2\quad 平拋運動}

\subsection*{\sffamily A.\ 什麼是平拋}

把物體\textbf{水平}拋出（初速只有水平分量、沒有鉛直分量），讓它在重力下落，這種運動稱為\textbf{平拋運動（horizontal projectile）}。例如桌邊滾出的球、水平射出的子彈、飛機水平投下的物資。

依運動獨立性，分兩個方向處理。取拋出點為原點、水平向右為 $x$ 正向、鉛直向下為 $y$ 正向，初速為 $v_0$（水平）：

\begin{center}
\renewcommand{\arraystretch}{1.3}
\begin{tabular}{c c c c}
\hline
\textbf{方向} & \textbf{加速度} & \textbf{初速} & \textbf{運動類型} \\
\hline
水平 $x$ & $a_x=0$ & $v_0$ & 等速度 \\
鉛直 $y$ & $a_y=g$ & $0$ & 自由落體（等加速度） \\
\hline
\end{tabular}
\end{center}

對應的位置與速度方程式為
$$\boxed{\,x=v_0\,t,\qquad y=\tfrac12 g t^2\,}$$
$$v_x=v_0\ (\text{不變}),\qquad v_y=g\,t,\qquad v=\sqrt{v_0^2+(gt)^2}.$$
水平分量 $v_x$ 從頭到尾都等於 $v_0$，不會減少（沒有水平方向的力）；鉛直分量 $v_y$ 像自由落體一樣越來越大。因此合速度越來越偏向下方。

\fig{si3-projectile}{0.74}{圖 3-2：平拋的速度分解。水平分量 $v_x=v_0$ 保持不變，鉛直分量 $v_y=gt$ 越來越大，合速度（紫色）越來越偏向下方。}

\subsection*{\sffamily B.\ 軌跡是拋物線}

由 $x=v_0 t$ 解出 $t=\dfrac{x}{v_0}$，代入 $y=\tfrac12 g t^2$：
$$y=\frac12 g\left(\frac{x}{v_0}\right)^2=\frac{g}{2v_0^2}\,x^2.$$
這是 $y=k x^2$ 的形式（$k$ 為常數），即一條標準的拋物線。

\begin{補充框}{拋體的軌跡為什麼一定是拋物線？}
「拋物線（parabola）」是數學上有嚴格定義的曲線。拋體之所以剛好走這條曲線，是\textbf{推導}出來的結果，而非直接宣布。

關鍵在於\textbf{水平位置與時間成正比、鉛直位置與時間平方成正比}。把時間消去後，鉛直位置就與水平位置的平方成正比，而「$y$ 正比於 $x^2$」正是拋物線的定義式。以平拋為例，由 $x=v_0 t$ 得 $t=\dfrac{x}{v_0}$，代入 $y=\tfrac12 g t^2$ 即得 $y=\dfrac{g}{2v_0^2}x^2$。

斜拋也相同，只是式子多一項。由 $x=(v_0\cos\theta)t$ 消去 $t$，代入 $y=(v_0\sin\theta)t-\tfrac12 g t^2$：
$$y=(\tan\theta)\,x-\frac{g}{2v_0^2\cos^2\theta}\,x^2,$$
仍是「一次項減去 $x^2$ 項」的形式，為開口向下的拋物線，只是頂點不在原點、整條被斜放。簡言之，「等速度 ＋ 等加速度」的疊加，在數學上必然是拋物線。

\textbf{注意：}拋體軌跡\textbf{不是}圓弧。圓上每一點到圓心等距，拋物線沒有這個性質，兩者只是局部形狀相近。
\end{補充框}

\subsection*{\sffamily C.\ 落地時間與水平射程}

從高 $H$ 處平拋：
\begin{itemize}\setlength{\itemsep}{1pt}
\item \textbf{落地時間} $t$ 由鉛直方向決定：$H=\tfrac12 g t^2\Rightarrow t=\sqrt{\dfrac{2H}{g}}$。此時間與水平初速 $v_0$ 無關。
\item \textbf{水平射程}（落地時的水平距離）：$R=v_0\,t=v_0\sqrt{\dfrac{2H}{g}}$。
\end{itemize}

\textbf{注意：}
\begin{itemize}\setlength{\itemsep}{1pt}
\item 「拋得越快，飛在空中越久」並不正確。落地時間只由高度決定；拋得快只是飛得遠，不會飛得久。
\item 「水平速度會慢慢變小」並不正確。忽略空氣阻力時 $v_x=v_0$ 永遠不變，變的是鉛直分量。
\item 落地速率不是 $v_0$，須把 $v_x$、$v_y$ 合起來：$v=\sqrt{v_0^2+(gt)^2}$。
\end{itemize}

\begin{補充框}{平拋的球與直接放手落下的球，哪個先落地？}
答案是\textbf{同時落地}。「多久落地」只由鉛直方向決定，而兩顆球的鉛直運動完全相同——都從靜止開始、以同一個 $g$ 自由下墜。水平方向跑多快，與「落地要多久」無關。

只看鉛直方向：兩顆球的鉛直初速都是 $0$（平拋的初速為純水平，鉛直分量為零）、鉛直加速度都是 $g$、起始高度都是 $H$，故落地時間都是 $t=\sqrt{\dfrac{2H}{g}}$。平拋的球只是額外往水平方向跑了 $R=v_0 t$，這段水平位移不占用鉛直方向的時間。

\textbf{注意：}「拋出去的球走斜線、路徑較長，所以較慢落地」並不正確。路徑長度與落地時間沒有直接關係，決定時間的是鉛直方向的運動。至於\textbf{落地速率}，平拋球的落地速率 $\sqrt{v_0^2+(gt)^2}$ 大於純落下的 $gt$；「較快」指的是速率，不是「較早到」。
\end{補充框}

\fig{si3-drop-compare}{0.50}{圖 3-3：水平拋出與直接放手的兩顆球，每隔相同時間並排畫出。它們在每一個時刻的高度都相同（灰色虛線等高），因此同時落地。}

\section*{\sffamily 3-3\quad 斜向拋射}

\subsection*{\sffamily A.\ 斜拋的初速分解}

把物體以與水平成 $\theta$ 角、初速大小 $v_0$ 斜向上拋出，稱為\textbf{斜向拋射（斜拋，oblique projectile）}。處理的第一步是把初速分解成水平、鉛直兩個分量：
$$v_{0x}=v_0\cos\theta\ (\text{水平}),\qquad v_{0y}=v_0\sin\theta\ (\text{鉛直，向上}).$$

\fig{si3-oblique}{0.92}{圖 3-4：左——斜拋初速分解為水平 $v_0\cos\theta$ 與鉛直 $v_0\sin\theta$。右——相同初速、不同角度的軌跡；$45^\circ$ 射程最大，互補的兩角（如 $30^\circ$ 與 $60^\circ$）落點相同。}

接著各方向獨立處理（取向上為正，$g$ 向下）：

\begin{center}
\renewcommand{\arraystretch}{1.3}
\begin{tabular}{c c c c}
\hline
\textbf{方向} & \textbf{加速度} & \textbf{初速} & \textbf{運動類型} \\
\hline
水平 $x$ & $0$ & $v_0\cos\theta$ & 等速度 \\
鉛直 $y$ & $-g$ & $v_0\sin\theta$ & 鉛直上拋（等加速度） \\
\hline
\end{tabular}
\end{center}

對應的位置與速度方程式為
$$x=(v_0\cos\theta)\,t,\qquad y=(v_0\sin\theta)\,t-\tfrac12 g t^2,$$
$$v_x=v_0\cos\theta\ (\text{不變}),\qquad v_y=v_0\sin\theta-g\,t.$$
斜拋即「水平等速度 ＋ 鉛直上拋」的疊加。鉛直方向與鉛直上拋相同，因此「最高點 $v_y=0$」「上升時間等於下降時間」等結論都能直接沿用。

\subsection*{\sffamily B.\ 飛行時間與最高高度}

\textbf{最高點}為鉛直速度為零的瞬間，$v_y=0$：
$$v_0\sin\theta-g\,t_{\text{頂}}=0\Rightarrow t_{\text{頂}}=\frac{v_0\sin\theta}{g}.$$

\textbf{總飛行時間}（回到同一水平高度，上升與下降對稱）：
$$\boxed{\,T=2t_{\text{頂}}=\frac{2v_0\sin\theta}{g}\,}$$

\textbf{最高高度} $H$（由鉛直方向 $v_y^2=v_{0y}^2-2gH$，頂點 $v_y=0$）：
$$\boxed{\,H=\frac{v_0^2\sin^2\theta}{2g}\,}$$

\subsection*{\sffamily C.\ 水平射程公式}

\textbf{水平射程} $R$＝水平速度 $\times$ 總飛行時間：
$$R=(v_0\cos\theta)\cdot T=(v_0\cos\theta)\cdot\frac{2v_0\sin\theta}{g}=\frac{2v_0^2\sin\theta\cos\theta}{g}.$$
用二倍角公式 $2\sin\theta\cos\theta=\sin 2\theta$：
$$\boxed{\,R=\frac{v_0^2\sin 2\theta}{g}\,}$$

由此得到兩個推論：
\begin{itemize}\setlength{\itemsep}{1pt}
\item \textbf{$45^\circ$ 射程最大}：$\sin 2\theta$ 在 $2\theta=90^\circ$（即 $\theta=45^\circ$）時最大（$=1$），此時 $R_{\max}=\dfrac{v_0^2}{g}$。
\item \textbf{互補角射程相同}：$\theta$ 與 $90^\circ-\theta$ 的 $\sin 2\theta$ 相等（如 $30^\circ$ 與 $60^\circ$），故射程相同（但飛行時間、最高高度不同）。
\end{itemize}

\textbf{注意：}射程公式 $R=\dfrac{v_0^2\sin 2\theta}{g}$ 只在\textbf{落回原拋出高度}時成立。若落點低於或高於拋出點（如從崖上斜拋落到崖底），不能直接套用，須回到 $x$、$y$ 方程式求解（見 3-4）。另外，最高點時並非「速度為零」，而是只有鉛直分量 $v_y=0$；水平分量 $v_x=v_0\cos\theta$ 仍存在，故最高點的速率為 $v_0\cos\theta$。

\begin{延伸框}{$45^\circ$ 射程最大的物理原因；有空氣阻力時是否成立}
公式上「$45^\circ$ 最遠」來自 $\sin 2\theta$ 在 $45^\circ$ 達到最大值，其背後的物理是一場\textbf{「水平要快」與「滯空要久」的拉鋸}。把射程寫成
$$R=\underbrace{(v_0\cos\theta)}_{\text{水平速度}}\times\underbrace{\frac{2v_0\sin\theta}{g}}_{\text{滯空時間}}.$$
角度越小，水平速度越大，但鉛直初速越小、滯空時間越短；角度越大則相反。兩者乘積的最佳妥協點落在 $45^\circ$，此時固定的速率平均分配給水平與鉛直分量（$\cos45^\circ=\sin45^\circ$）。也可由對稱性直接看出：射程公式對 $\theta$ 與 $90^\circ-\theta$ 對稱，對稱函數的極值落在對稱中心 $\dfrac{0^\circ+90^\circ}{2}=45^\circ$。

\textbf{有空氣阻力時，最佳角通常小於 $45^\circ$。}空氣阻力大致與速率有關、方向與運動方向相反；飛得越久、越高，被阻力消耗的能量越多。為減少阻力影響，較有利的方式是較低、較早地飛出，因此最佳角會降到 $45^\circ$ 以下，且阻力越大、初速越快，降得越多。實測的最佳發射角大致為：鉛球約 $42^\circ$（重、慢，阻力影響小），標槍、鏈球約 $30^\circ\!\sim\!40^\circ$，高爾夫球、棒球長打常更低，並牽涉旋轉產生的升力（馬格努斯效應）。

\textbf{此外，出手高度也有影響。}即使沒有空氣阻力，若落點低於出手點，射程公式因「落回同高度」的前提不成立而不適用；此時最佳角同樣小於 $45^\circ$，因為站在高處相當於先賺得一段滯空時間。「$45^\circ$ 最遠」是\textbf{無阻力、落回同高度}的理想模型結論，標清前提即無矛盾。
\end{延伸框}

\section*{\sffamily 3-4\quad 拋體運動的綜合分析}

前三節為基本題型。當條件改變、現成公式無法直接套用時，回到最原始的兩條方程式求解：
$$x=(v_0\cos\theta)\,t,\qquad y=(v_0\sin\theta)\,t-\tfrac12 g t^2.$$
把已知代入、解出時間 $t$，其餘各量隨之而得。例如從高處斜拋、落到較低地面時，落地點的 $y$ 為負值（在拋出點下方），代入鉛直方程式解出 $t$ 後，再由水平方程式得到水平距離；此情形不能使用射程公式 $R=\dfrac{v_0^2\sin2\theta}{g}$。

\textbf{注意：}拋體在某一時刻的速度方向，是\textbf{合速度的方向}（沿軌跡切線），由 $\tan\alpha=\dfrac{v_y}{v_x}$ 決定，會隨時間改變；它與「位移方向」「初速方向」不同。最高點時 $v_y=0$，速度方向恰為水平。

\medskip
運動獨立性的一個直接結論是：兩個物體若在鉛直方向受到相同的重力，則它們的鉛直運動相同。據此，一顆\textbf{水平射出的子彈}與一個\textbf{同時自由落下的目標}，在任一時刻下墜的高度 $\tfrac12 g t^2$ 完全相同。若忽略重力，子彈會沿直線飛向目標原來的位置；加入重力後，子彈相對於該直線下墜多少，目標也正好下墜多少。因此只要子彈在目標落地前抵達其水平位置，兩者必在空中相遇。這說明子彈的鉛直下墜與其水平速度互不影響，是運動獨立性的典型應用。

\section*{\sffamily 3-5\quad 平面（二維）相對運動}

\subsection*{\sffamily A.\ 從一維相對速度升級為向量相對速度}

直線運動中已介紹\textbf{一維相對速度}：在同一條直線上，B 相對於 A 的速度為兩者相減，$v_{B/A}=v_B-v_A$。但「在運動中觀察另一物體」不只發生在一條線上：在河上行船、在風中飛行時，水流與風的方向常和欲前進的方向\textbf{垂直}。此時須把相對速度當成\textbf{向量}處理，等於把一維的想法升級為兩個分量的向量版。式子與一維時完全相同，只是把純量換成向量：
$$\boxed{\;\vec v_{B/A}=\vec v_B-\vec v_A\;}$$
讀作「B 相對於 A 的速度＝B 對地速度 $-$ A 對地速度」，意即站在 A 上觀察 B 所看到的速度。差別只在這裡是\textbf{向量相減}——須分 $x$、$y$ 兩個分量各自相減，最後再合成。

\begin{補充框}{對地速度＝對介質速度＋介質對地速度}
把上式移項，可得實務上最常用的關係式。以船為例（取 A＝水、B＝船）：
$$\underbrace{\vec v_{\text{船對地}}}_{\text{合速度（真正航跡）}}=\underbrace{\vec v_{\text{船對水}}}_{\text{划船的努力}}+\underbrace{\vec v_{\text{水對地}}}_{\text{水流}}.$$
船真正的航跡（對地速度）是「船自己划的速度」與「水流速度」兩個向量\textbf{頭尾相接}合成的結果。把這兩個已知向量頭尾相接、第三邊即為合速度，便得到一個\textbf{向量三角形}；解題時務必標清楚哪一邊是「想要前進的方向」。

風中飛機是同一個模型，只是把「水」換成「空氣」：對地速度（航跡）＝對空速度（機頭朝向、空速）＋風速。甚至「在傳送帶上行走」也是同一件事，皆可歸結為「對地＝對介質＋介質對地」這一張向量三角形。
\end{補充框}

\subsection*{\sffamily B.\ 河中行船：兩種典型問法}

設河寬 $d$、水流速度大小 $u$（沿河岸方向）、船相對水的速度大小 $v$。因為水流（沿河岸方向）與船划向對岸（垂直河岸方向）互不干涉——這又是一次運動獨立性——所以兩個方向可分開計算。取\textbf{垂直河岸}為 $y$、\textbf{沿河岸（順流）}為 $x$：只有船速在 $y$ 方向的分量帶著船過河，故\textbf{渡河時間只由這個分量決定，與水流 $u$ 無關}；$x$ 方向則被水流推動，造成往下游的\textbf{漂移}。

\textbf{問法一：船頭垂直指向對岸（最快渡河）。} 船把全部努力都用在過河方向，$y$ 方向分量最大（$=v$），故\textbf{渡河時間最短}：
$$t=\frac{d}{v}.$$
但水流整段時間都在推，船會被沖到下游 $x=u\,t=\dfrac{ud}{v}$ 處。此時對地的合速度大小為 $\sqrt{v^2+u^2}$，方向斜指下游。

\textbf{問法二：要垂直抵達正對岸（不被沖走）。} 須讓船頭朝\textbf{上游}偏一角度 $\theta$，使船速的「沿河岸分量」$v\sin\theta$ 恰好抵消水流 $u$：
$$v\sin\theta=u\;\Rightarrow\;\sin\theta=\frac{u}{v}.$$
此時真正過河的有效速度只剩 $v\cos\theta=\sqrt{v^2-u^2}$，渡河時間為
$$t=\frac{d}{\sqrt{v^2-u^2}},$$
比問法一久。顯然須 $v>u$ 才可能垂直渡河；若 $v\le u$，水流過強，無論如何划都會被沖向下游。

\textbf{注意：}
\begin{itemize}\setlength{\itemsep}{1pt}
\item 「水流會讓渡河變慢」並不一定成立。若船頭直指對岸（問法一），渡河時間 $t=\dfrac{d}{v}$ 與水流 $u$ 完全無關，水流只決定被沖多遠，不決定多久過河。會變慢的是「要垂直抵達正對岸」這種把部分船速拿去抵消水流的情形。
\item 「合速度大小是 $v+u$」只有兩者同向時才成立。船速與水流垂直時，合速度為 $\sqrt{v^2+u^2}$；一般情形須按向量（分量）合成，不能直接相加。
\end{itemize}

\subsection*{\sffamily C.\ 風中飛機：同一招}

飛機問題與行船是同一個模型，只是把水換成空氣：對地速度（航跡）＝對空速度（機頭朝向、空速）＋風速。常見問法是「機頭該朝哪個方向、實際對地飛多快」。解法完全相同：若有側風使航跡偏離目標方向，就讓機頭朝\textbf{側風來向}偏一角度 $\theta$，使空速的「側向分量」$v\sin\theta$ 抵消風速 $u$（即 $\sin\theta=\dfrac{u}{v}$）；沿目標方向真正前進的對地速率則為 $v\cos\theta=\sqrt{v^2-u^2}$。整個過程與「垂直抵達對岸」的行船問法是同一張向量三角形。

\section*{\sffamily 3-6\quad 本章重點整理}

\begin{補充框}{一頁複習（公式與關鍵結論）}
\begin{itemize}\setlength{\itemsep}{2pt}
\item \textbf{運動獨立性}：水平、鉛直兩方向互不影響，分開用一維公式計算，再合成。根源是重力只有鉛直分量。前提為忽略空氣阻力。
\item \textbf{平拋}（初速水平 $v_0$）：水平等速度 $x=v_0 t$、$v_x=v_0$；鉛直自由落體 $y=\tfrac12 g t^2$、$v_y=gt$。
  \begin{itemize}\setlength{\itemsep}{1pt}
  \item 從高 $H$ 落地：$t=\sqrt{\dfrac{2H}{g}}$（與 $v_0$ 無關）；射程 $R=v_0\sqrt{\dfrac{2H}{g}}$。
  \item 軌跡 $y=\dfrac{g}{2v_0^2}x^2$（拋物線）。
  \end{itemize}
\item \textbf{斜拋}（初速 $v_0$、仰角 $\theta$）：先分解 $v_{0x}=v_0\cos\theta$、$v_{0y}=v_0\sin\theta$。
  \begin{itemize}\setlength{\itemsep}{1pt}
  \item 飛行時間 $T=\dfrac{2v_0\sin\theta}{g}$；最高高度 $H=\dfrac{v_0^2\sin^2\theta}{2g}$；射程 $R=\dfrac{v_0^2\sin 2\theta}{g}$。
  \item $45^\circ$ 射程最大 $R_{\max}=\dfrac{v_0^2}{g}$；互補角（如 $30^\circ$ 與 $60^\circ$）射程相同。
  \end{itemize}
\item \textbf{通用法}（公式套不上時）：回到 $x=(v_0\cos\theta)t$、$y=(v_0\sin\theta)t-\tfrac12 gt^2$，解出 $t$。
\item 落地速率為合速度 $v=\sqrt{v_x^2+v_y^2}$；最高點 $v_y=0$ 但 $v_x\ne0$。
\item \textbf{平面相對速度}：$\vec v_{B/A}=\vec v_B-\vec v_A$（向量相減，分量各自算）；對地＝對介質＋介質對地。
  \begin{itemize}\setlength{\itemsep}{1pt}
  \item 行船／飛機：船頭直指對岸時\textbf{渡河時間 $t=\dfrac{d}{v}$ 與水流無關}，被沖到下游 $ut$ 處。
  \item 要垂直抵達對岸須朝上游偏 $\sin\theta=\dfrac{u}{v}$，渡河時間 $t=\dfrac{d}{\sqrt{v^2-u^2}}$（需 $v>u$）。
  \end{itemize}
\end{itemize}
\end{補充框}

\end{document}
